\chapter{Реализация прототипа}\label{ch:prototype}

В данном разделе описана реализация прототипа системы интеллектуальных автономных агентов: технологический стек, архитектура модулей, интеграция с LLM через Model Context Protocol, сценарии самовосстановления и результаты тестирования.

\section{Цели и область применения прототипа}

Прототип реализует архитектуру системы интеллектуальных агентов (раздел~\ref{ch:architecture}) в виде работающего программного комплекса. Цель прототипа --- проверить гипотезу~(\ref{eq:hypothesis_ineq}): обеспечить уменьшение среднего времени восстановления~$\text{MTTR}$ и улучшение метрик качества (точность, ложные срабатывания, время диагностики) на измеримые величины~$I_m$, определённые во введении.

Основные цели:
\begin{itemize}
  \item интегрировать модули сбора метрик, логов и трейсов с промышленными системами наблюдаемости (Prometheus, Elasticsearch, Jaeger);
  \item реализовать цепочку Anomaly Detector~$\to$ Diagnostic Engine~$\to$ \linebreak Orchestrator~$\to$ Safe Executor;
  \item обеспечить воспроизведение типовых сценариев инцидентов в условиях, близких к production;
  \item подготовить экспериментальную платформу для количественной оценки (раздел~\ref{ch:experiments}).

\end{itemize}

\section{Технологический стек}

Прототип реализован на Java с использованием Spring Boot \cite{spring_framework}. Технологический стек прототипа представлен в таблице~\ref{tab:tech_stack}. Выбор DeepLearning4j обусловлен требованием единой JVM-\linebreak экосистемы: все модули агента (Collector, Detector, Diagnostic, Executor) выполняются в JVM, что исключает межпроцессную сериализацию при передаче данных между детектором аномалий и остальными компонентами. Альтернативы (PyTorch, TensorFlow) потребовали бы Python-сервиса с REST/gRPC-мостом, увеличивая латентность инференса на 15--30 мс.

\begin{table}[H]
\centering
\caption{Технологический стек прототипа}
\label{tab:tech_stack}
\begin{tabular}{|p{3.5cm}|p{3.5cm}|p{7cm}|}
\hline
\textbf{Компонент} & \textbf{Технология} & \textbf{Назначение} \\
\hline
Язык и фреймворк & Java, Spring Boot & Основная платформа разработки \\
\hline
ML-инференс & DeepLearning4j & LSTM-модель обнаружения аномалий \\
\hline
Графовые алгоритмы & JGraphT & Построение причинно-следственных графов \\
\hline
Реактивное взаимодействие & Spring WebFlux & Асинхронная обработка потоков данных \\
\hline
Шина событий & Apache Kafka & Передача событий между модулями \\
\hline
Кэш контекста & Redis & Хранение оперативного контекста агентов \\
\hline
Хранение метрик & InfluxDB & Временные ряды метрик наблюдаемости \\
\hline
Управление кластером & Kubernetes Client & Выполнение действий восстановления \\
\hline
\end{tabular}
\end{table}

Кластер развёртывания включает:
\begin{itemize}
  \item Kubernetes-кластер с 10 сервисами бизнес-логики (микросервисное приложение e-commerce);
  \item подсистему наблюдаемости: Prometheus, Elasticsearch, Jaeger;
  \item подсистему агентов: Collector, Detector, Diagnostic, Orchestrator, Executor, LLM Agent --- каждый как отдельный \texttt{Deployment}.

\end{itemize}

Архитектура прототипа представлена на рисунке~\ref{fig:prototype_architecture}. Взаимодействие компонентов организовано через:
\begin{itemize}
  \item Kafka-топики для передачи событий аномалий и диагнозов;
  \item Redis как кэш контекстной информации;
  \item REST/gRPC-интерфейсы для вызова действий восстановления.

\end{itemize}

\begin{figure}[H]
\centering
\begin{tikzpicture}[node distance=0.8cm, auto, scale=0.75, transform shape, font=\small]
    \node[rectangle, draw, fill=blue!20, minimum width=1.2cm, minimum height=0.6cm] (prometheus) {Prometheus};
    \node[rectangle, draw, fill=green!20, below=of prometheus, minimum width=1.2cm, minimum height=0.6cm] (elasticsearch) {Elasticsearch};
    \node[rectangle, draw, fill=yellow!20, below=of elasticsearch, minimum width=1.2cm, minimum height=0.6cm] (jaeger) {Jaeger};
    \node[rectangle, draw, fill=red!20, right=of prometheus, xshift=1cm, minimum width=1.2cm, minimum height=0.6cm] (collector) {Collector};
    \node[rectangle, draw, fill=orange!20, right=of collector, xshift=1cm, minimum width=1.2cm, minimum height=0.6cm] (detector) {Detector};
    \node[rectangle, draw, fill=purple!20, below=of detector, minimum width=1.2cm, minimum height=0.6cm] (diagnostic) {Diagnostic};
    \node[rectangle, draw, fill=cyan!20, below=of collector, minimum width=1.2cm, minimum height=0.6cm] (orchestrator) {Orchestrator};
    \node[rectangle, draw, fill=pink!20, below=of orchestrator, minimum width=1.2cm, minimum height=0.6cm] (executor) {Executor};
    \node[rectangle, draw, fill=brown!20, right=of diagnostic, xshift=1cm, minimum width=1.2cm, minimum height=0.6cm] (llm) {LLM};

    \draw[->, shorten >=2pt, shorten <=2pt] (prometheus) -- (collector);
    \draw[->, shorten >=2pt, shorten <=2pt] (elasticsearch) -- (collector);
    \draw[->, shorten >=2pt, shorten <=2pt] (jaeger) -- (collector);
    \draw[->, shorten >=2pt, shorten <=2pt] (collector) -- (detector);
    \draw[->, shorten >=2pt, shorten <=2pt] (detector) -- (diagnostic);
    \draw[->, shorten >=2pt, shorten <=2pt] (diagnostic) -- (orchestrator);
    \draw[->, shorten >=2pt, shorten <=2pt] (orchestrator) -- (executor);
    \draw[->, shorten >=2pt, shorten <=2pt] (diagnostic) -- (llm);
    \draw[->, shorten >=2pt, shorten <=2pt] (llm) -- (orchestrator);
\end{tikzpicture}
\caption{Архитектура прототипа}
\label{fig:prototype_architecture}
\end{figure}

\section{Реализация модулей}

Каждый модуль прототипа соответствует компоненту архитектуры из раздела~\ref{ch:architecture}. Ниже описаны особенности реализации, не покрытые формальными моделями.

\subsection{Модуль сбора и нормализации данных}

Сервис \texttt{DataCollectorService} реализует периодический и событийный опрос Prometheus, Elasticsearch и Jaeger, приводя данные к унифицированному формату (\ref{eq:unified_data_model}). \texttt{DataProcessorService} выполняет агрегацию, фильтрацию и нормализацию: усреднение по окнам, downsampling, обогащение контекстом.

\paragraph{Пример унификации.}

В качестве примера рассмотрим метрику latency сервиса checkout. В Prometheus она представлена как \linebreak \texttt{http\_request\_duration\_seconds\{service="checkout"\}} со значением $0.8$~c. Elasticsearch содержит лог уровня ERROR с \texttt{trace\_id}~$= \tau$. Collector преобразует их в записи:
\begin{align*}
  d_1 &= (t, \text{``checkout''}, \text{metric}, 0.8, \\
      &\qquad \{\text{status} = \text{``5xx''}\}), \\
  d_2 &= (t, \text{``checkout''}, \text{log}, \text{``ERROR''}, \\
      &\qquad \{\text{trace\_id} = \tau\}),
\end{align*}
что позволяет связывать метрики и логи по полям \(\text{source}\) и \(\text{trace\_id}\) при построении диагностического графа.

\subsection{Модуль обнаружения аномалий}

\texttt{AnomalyDetectorService} реализует алгоритмы, описанные в разделе~\ref{ch:anomaly}: пороговые правила (Z-score), модели ARIMA и многомерную LSTM-модель. Конфигурация LSTM: $w = 60$ точек, $k = 3$ (параметр чувствительности), критерий аномалии $e_t > \mu + k\sigma$.

Результат работы --- события \linebreak \texttt{AnomalyEvent(service, metric, score, window)}, размещаемые в Kafka-топиках. Например, при $Z(t) = 3.67 > \theta_Z = 3$ для метрики latency сервиса \texttt{checkout} и одновременном LSTM-прогнозе $\hat{\text{Latency}}(t+\Delta t) = 380 > \theta_{\text{degradation}} = 350$ мс формируется событие аномалии, передаваемое в Diagnostic Engine.

\subsection{Модуль диагностического графа}

\texttt{DiagnosticEngineService} строит причинно-следственный граф по формулам~(\ref{eq:causal_probability})--(\ref{eq:root_cause}) с использованием библиотеки JGraphT. На вход поступают метрики, логи и трейсы; на выходе --- \texttt{DiagnosisReport} с ранжированными кандидатами на корневую причину.

\paragraph{Пример.}
Для инцидента с ростом latency \texttt{/checkout} трейс Jaeger фиксирует цепочку
$\text{API-Gateway} \rightarrow \text{CheckoutService} \rightarrow \text{PaymentService} \rightarrow \text{DB}$.
Модуль строит граф, где $P(\text{PaymentService} \to \text{DB}) \approx 0.8$ (высокая корреляция задержек и топология трейса), а $P(\text{CheckoutService} \to \text{PaymentService}) \approx 0.3$. С учётом ошибок 5xx только на уровне \texttt{PaymentService} алгоритм выдаёт $\text{RootCause} = \text{PaymentService}$.

\subsection{Модуль оркестрации и исполнения}

Модуль \texttt{orchestrator} реализует консенсус агентов по формулам~(\ref{eq:consensus})--(\ref{eq:degroot_consensus}). На вход поступают рекомендации от Diagnostic Engine, LLM-агента и rule-based политик; на выходе --- \texttt{RecoveryPlan}.

\texttt{RecoveryExecutorService} преобразует план в операции Kubernetes: масштабирование, изменение конфигурации, управление circuit breaker. Все действия выполняются через Kubernetes Client. Все действия валидируются модулем Safe Executor по критерию~(\ref{eq:safety_criterion}) с порогом $r_{\max} = 0.3$.

\subsection{Модуль LLM Agent}

LLM-агент получает структурированный контекст и формирует: текстовые объяснения инцидента, гипотезы причин и рекомендации по восстановлению. Результаты поступают как в оркестратор, так и в каналы разработчиков (Slack/почта). Интеграция реализована через Model Context Protocol (раздел~\ref{sec:mcp}).

\section{Интеграция с LLM через Model Context Protocol}
\label{sec:mcp}

MCP --- протокол на основе JSON-RPC~2.0 для стандартизированного взаимодействия LLM-агентов с внешними системами \cite{mcp_specification}. MCP-сервер выступает прослойкой между LLM-провайдерами и внутренними сервисами прототипа.

\subsection{Реализованные MCP-инструменты}

В прототипе реализовано пять инструментов, перечисленных в таблице~\ref{tab:mcp_tools}, каждый из которых соответствует логически целостной операции:

\begin{table}[H]
\centering
\caption{MCP-инструменты прототипа}
\label{tab:mcp_tools}
\begin{tabular}{|p{4cm}|p{7.5cm}|}
\hline
\textbf{Инструмент} & \textbf{Назначение} \\
\hline
\texttt{system\_context} & Агрегированный срез состояния: метрики $M(t)$, сервисы $S(t)$, события $E(t-\Delta t, t)$, алерты $A(t)$ \\
\hline
\texttt{log\_analysis} & Фильтрация ошибок, извлечение паттернов и сущностей (NER), корреляции логов с метриками \\
\hline
\texttt{causal\_graph} & Доступ к диагностическому графу: вершины, рёбра, кандидаты на корневую причину, вероятности \\
\hline
\texttt{recommendations} & Полный цикл: контекст + логи + граф $\to$ LLM $\to$ \linebreak резюме, причины, рекомендуемые действия \\
\hline
\texttt{action\_result} & Фиксация результата: действие принято/выполнено/откатано/отклонено \\
\hline
\end{tabular}
\end{table}

Схема взаимодействия LLM-провайдера с MCP-сервером показана на рисунке~\ref{fig:mcp_server}. Обработка каждого запроса реализована как композиция этапов:
\begin{equation}
\text{Process}(R) = \text{Validate} \to \text{Authenticate} \to \text{Execute} \to \text{Format}
\label{eq:mcp_processing}
\end{equation}

\begin{figure}[H]
\centering
\begin{tikzpicture}[node distance=1.5cm, auto, scale=0.75, transform shape, font=\small]
    \node[rectangle, draw, fill=blue!20, minimum width=2cm, minimum height=0.8cm] (llm) {LLM-провайдер};
    \node[rectangle, draw, fill=green!20, below=of llm, minimum width=2.5cm, minimum height=0.8cm] (mcp) {MCP-сервер};
    \node[rectangle, draw, fill=yellow!20, below left=of mcp, minimum width=2cm, minimum height=0.8cm] (ctx) {system\_context};
    \node[rectangle, draw, fill=yellow!20, below=of mcp, minimum width=2cm, minimum height=0.8cm] (logs) {log\_analysis};
    \node[rectangle, draw, fill=yellow!20, below right=of mcp, minimum width=2cm, minimum height=0.8cm] (graph) {causal\_graph};

    \draw[->, shorten >=2pt, shorten <=2pt] (llm) -- node[right]{JSON-RPC} (mcp);
    \draw[->, shorten >=2pt, shorten <=2pt] (mcp) -- (ctx);
    \draw[->, shorten >=2pt, shorten <=2pt] (mcp) -- (logs);
    \draw[->, shorten >=2pt, shorten <=2pt] (mcp) -- (graph);
    \draw[->, shorten >=2pt, shorten <=2pt] (ctx) -- (mcp);
    \draw[->, shorten >=2pt, shorten <=2pt] (logs) -- (mcp);
    \draw[->, shorten >=2pt, shorten <=2pt] (graph) -- (mcp);
    \draw[->, shorten >=2pt, shorten <=2pt] (mcp) -- node[left]{Ответ} (llm);
\end{tikzpicture}
\caption{Схема взаимодействия LLM-провайдера с MCP-сервером}
\label{fig:mcp_server}
\end{figure}

Типичный диалог: LLM формирует запрос <<объяснить рост 5xx для \texttt{order-service}>>, MCP-сервер вызывает \texttt{system\_context}, \texttt{log\_analysis} и \texttt{causal\_graph}, агрегирует результаты и возвращает структурированный ответ. Поддерживаются провайдеры: OpenAI (GPT-4), Anthropic (Claude~3), Ollama (Llama~3), Azure OpenAI.

\subsection{Отказоустойчивость MCP-сервера}

Retry с экспоненциальной задержкой:
\begin{equation}
T_{\text{retry}} = T_{\text{base}} \cdot 2^{n-1} \cdot (1 + \text{random}(0, j))
\label{eq:retry_delay}
\end{equation}

Кэширование контекста с TTL, определяемым минимальным интервалом обновления источников:
\begin{equation}
\tau_{\text{cache}} = \min(\Delta t_{\text{metrics}}, \Delta t_{\text{logs}}, \Delta t_{\text{traces}})
\label{eq:cache_ttl}
\end{equation}

\section{Сценарии самовосстановления}

Прототип реализует четыре сценария, моделирующих типовые инциденты в микросервисных системах. Общий цикл самовосстановления, реализованный в прототипе, формализован в алгоритме~\ref{alg:recovery_cycle}. Динамика метрик при самовосстановлении для сценария~1 показана на рисунке~\ref{fig:recovery_scenario}.

\begin{algorithm}[H]
\caption{Общий цикл самовосстановления}
\label{alg:recovery_cycle}
\begin{algorithmic}[1]
\REQUIRE Поток данных наблюдаемости $D = \{M, L, T\}$
\ENSURE Восстановленное состояние системы
\STATE $score \gets \text{AnomalyDetector}(D)$
\IF{$score > \theta_{\text{anomaly}}$}
  \STATE $G \gets \text{BuildCausalGraph}(D)$
  \STATE $root \gets \text{FindRootCause}(G)$
  \STATE $action \gets \text{SelectAction}(root)$
  \IF{$\text{Risk}(action) < r_{\max}$}
    \STATE $\text{Execute}(action)$
    \STATE $\text{Verify}(action)$
  \ELSE
    \STATE $\text{Escalate}(action)$
  \ENDIF
\ENDIF
\end{algorithmic}
\end{algorithm}

\subsection{Сценарий 1: Отказ одного сервиса}

Моделируется внезапное завершение процесса Order Service (kill~-9). Обнаружение: исчезновение метрик и ошибки в трейсах зависимых сервисов. Диагностика: Diagnostic Engine строит причинно-следственный граф и идентифицирует отказавший сервис как корневую причину. Действие --- автоматический перезапуск:
\begin{equation}
N_{\text{new}} = \max\!\left(N_{\text{desired}},\; N_{\text{current}} + 1\right)
\label{eq:scaling_action}
\end{equation}
Валидация: восстановление health check и нормализация error rate зависимых сервисов.

\subsection{Сценарий 2: Каскадный отказ}

Моделируется перегрузка Payment Service (CPU stress 95\,\%), вызывающая таймауты в Order Service и Notification Service. Цепочка зависимостей: Payment $\to$ Order $\to$ API~Gateway. Обнаружение: рост $\text{ErrorRate}_{5xx}$ и latency на нескольких сервисах одновременно. Диагностика: локализация корневой причины по причинно-следственному графу --- идентификация Payment Service как источника каскада.

Действие --- комплексное восстановление с оценкой ожидаемого времени каскадного распространения:
\begin{equation}
T_{\text{cascade}} = \sum_{i=1}^{n} \Delta t_i \cdot P(s_i \mid s_{i-1}), \quad \Delta t_i = t_{\text{fail}}(s_i) - t_{\text{fail}}(s_{i-1})
\label{eq:cascade_time}
\end{equation}
где $n$ --- глубина цепочки зависимостей, $\Delta t_i$ --- интервал между отказами смежных сервисов $s_{i-1}$ и $s_i$, $P(s_i \mid s_{i-1})$ --- условная вероятность отказа $s_i$ при отказе $s_{i-1}$. На основе $T_{\text{cascade}}$ оркестратор определяет порядок восстановления: масштабирование Payment Service, затем включение circuit breaker на зависимых сервисах (Order, API~Gateway). Валидация: возврат $\text{ErrorRate}$ и latency к допустимым уровням по всей цепочке.

\subsection{Сценарий 3: Утечка памяти}

Моделируется постепенная деградация Catalog Service (memory leak 50~МБ/мин), приводящая к OOM-kill через 30--60 минут. Обнаружение: LSTM-детектор фиксирует аномальный тренд потребления памяти и прогнозирует время до OOM:
\begin{equation}
T_{\text{OOM}} = \arg\min_t \{t : \hat{M}_{\text{mem}}(t) > \theta_{\text{OOM}}\}
\label{eq:degradation_prediction_proto}
\end{equation}
Действие: превентивный перезапуск сервиса до наступления OOM, увеличение реплик для сохранения доступности. Сценарий демонстрирует использование LSTM не только для фиксации факта деградации, но и для упреждающего самовосстановления.

\subsection{Сценарий 4: Сетевая деградация}

Моделируется увеличение латентности сети между сервисом и базой данных (задержка 200~мс, packet loss 5\%). Проблема проявляется через рост таймаутов и ошибок 504. Diagnostic Engine строит причинно-следственный граф, выявляя сетевой сегмент как корневую причину. Safe Executor инициирует перенаправление трафика на резервный путь и ограничение нагрузки по критерию~(\ref{eq:safety_criterion}).

\begin{figure}[H]
\centering
\begin{tikzpicture}[scale=0.65]
\begin{axis}[
    xlabel={Время, сек},
    ylabel={CPU, \%},
    title={Сценарий самовосстановления при отказе сервиса},
    grid=major,
    legend pos=north west,
    width=0.9\textwidth,
    height=6cm,
]
\addplot[blue,thick] coordinates {
    (0, 50) (10, 60) (20, 75) (30, 92) (40, 95) (50, 65) (60, 55)
};
\addplot[red,dashed,thick] coordinates {
    (30, 95) (40, 65)
};
\node at (axis cs:35,80) {\small Восст.};
\legend{Загрузка CPU, Восстановление}
\end{axis}
\end{tikzpicture}
\caption{Динамика загрузки CPU при отказе сервиса (сценарий~1)}
\label{fig:recovery_scenario}
\end{figure}

\section{Конфигурация и развёртывание}

Параметры системы задаются в \texttt{application.yml} модулей и манифестах Kubernetes. Ключевые параметры перечислены в таблице~\ref{tab:prototype_params}.

\begin{table}[H]
\centering
\caption{Основные параметры конфигурации прототипа}
\label{tab:prototype_params}
\begin{tabular}{|p{5cm}|c|p{5cm}|}
\hline
\textbf{Параметр} & \textbf{Значение} & \textbf{Обоснование} \\
\hline
Порог аномалии ($\theta_{\text{anomaly}}$) & 0.85 & Grid search по $F_1$ (раздел~\ref{ch:experiments}) \\
\hline
Окно LSTM ($w_{\text{LSTM}}$) & 60 & Баланс детализации и вычислительных затрат \\
\hline
Максимум действий за цикл & 5 & Ограничение каскадных изменений \\
\hline
Порог риска ($r_{\max}$) & 0.3 & Экспертная оценка, подтверждена в экспериментах \\
\hline
Порог сходимости консенсуса ($\epsilon$) & 0.01 & Компромисс скорости и точности \\
\hline
\end{tabular}
\end{table}

Деплой запускается через манифесты Kubernetes с поэтапным обновлением и проверкой health checks. Для каждого сервиса агентов определены профили окружений (локальное развёртывание, тестовый кластер).

\section{Производительность и масштабируемость}

Измеренная производительность прототипа представлена в таблице~\ref{tab:prototype_performance}.

\begin{table}[H]
\centering
\caption{Производительность прототипа}
\label{tab:prototype_performance}
\begin{tabular}{|p{6cm}|c|}
\hline
\textbf{Показатель} & \textbf{Значение} \\
\hline
Обработка метрик (DataCollectorService) & $10^5$ метрик/с \\
\hline
Индексация логов (Elasticsearch) & $10^6$ сообщений/мин \\
\hline
Построение графа ($10^4$ вершин) & $\leq 5$ с \\
\hline
LLM-запросы через MCP & до 100 запросов/мин \\
\hline
Задержка Kafka (сбор $\to$ обнаружение) & $< 100$ мс \\
\hline
\end{tabular}
\end{table}

Масштабирование: для 10--50 сервисов --- один экземпляр каждого агента; для 50--500 --- 3--5 экземпляров с балансировкой по Kafka-топикам; для 500+ --- шардирование по ключам \texttt{service}/\texttt{namespace}.

Дискретизация: $\Delta t = 1$ с для метрик, $\Delta t = 100$ мс для трейсов; ошибка аппроксимации $\epsilon < 0.01$. Модульная декомпозиция проверена на отсутствие циклических зависимостей через статический анализ графа Maven.

\section{Тестирование}

Качество реализации подтверждено тремя уровнями тестирования:
\begin{itemize}
  \item \textbf{Unit-тесты} (JUnit~5): покрытие основных сервисов $> 80\%$;
  \item \textbf{Интеграционные тесты}: взаимодействие модулей через Kafka и Redis --- корректность форматов, обработка отказов брокера, повторные попытки;
  \item \textbf{End-to-end тесты}: полный цикл на реальных сценариях инцидентов (сбор данных $\rightarrow$ обнаружение $\rightarrow$ диагностика $\rightarrow$ восстановление $\rightarrow$ верификация).
\end{itemize}

Листинги ключевых модулей приведены в Приложении~В.

\section{Ограничения прототипа}

Ограничения прототипа:

\begin{itemize}
    \item \textbf{Целевая платформа.} Система ориентирована на микросервисные архитектуры с HTTP/gRPC-коммуникациями в Kubernetes. Инциденты на уровне сетевого оборудования (L2/L3), мэйнфреймов и serverless-функций не рассматриваются;

    \item \textbf{Набор действий восстановления.} Реализовано 4 типа действий: масштабирование реплик, перезапуск подов, изменение ConfigMap, управление circuit breaker. Миграция данных, переконфигурация топологии и управление DNS не автоматизированы;

    \item \textbf{Зависимость от LLM-провайдера.} Качество рекомендаций LLM Agent зависит от выбранной модели и её доступности. При недоступности LLM-провайдера (сетевой сбой, rate limit) система продолжает работать без LLM-обогащения, но с потерей 4.6\% $F_1$ (ablation study, раздел~\ref{ch:experiments}). Стоимость LLM-инференса (0.01--0.03 USD за запрос для GPT-4) не включена в TCO (раздел~\ref{ch:economic});

    \item \textbf{Качество входных данных.} LSTM-детектор калиброван на данных с интервалом scrape 15~с; при увеличении интервала (30~с, 60~с) пороги требуют рекалибровки. Downsampling метрик до 1-минутного разрешения снижает $F_1$ обнаружения на $\sim$8\%.
\end{itemize}
